\subsection{Composition of domain-changing interpretations}
\label{sec:p02-translation-composition}

A symbol interpretation and the operation it induces on a model must be
distinguished. Renaming symbols preserves the ambient universe. A
relativization can discard objects and can make a constant undefined.
These effects are not recorded by componentwise composition of symbol
maps alone.

\begin{hbtheorem}{thm:p02-translation-composition}{proved=annex:p2.tr.1}{The composition obstruction}
On the full category of finite languages and interpretations described
in \cref{def:p02-w02c-p02-003-b74-languages}, the following three
requirements cannot all hold:
\begin{enumerate}
\item sentences and their translations form a functor on that whole
      category, with the Occurrence Axiom;
\item the guarded Translation Axiom holds on the least language of
      each sentence;
\item atomic sentences have their usual truth conditions.
\end{enumerate}
A counterexample uses two objects, a unary predicate, one constant,
and two domain-changing zero-term interpretations. In particular,
ordinary first-order syntax with its usual semantics does not supply
such a functor for all those interpretations.
\end{hbtheorem}

These are the literal simultaneous requirements in the source's
Sections 4 and 7--8
\cite[pp.~229--230, 234--238]{Barwise1974Axioms}. The obstruction concerns
their composition law, not the validity of a single guarded
relativization. The working translation interface below keeps
functoriality for universe-preserving symbol changes and treats a
domain-changing translation as an explicitly guarded operation.
Sequential guarded translations remain available; they are not silently
identified with the componentwise composite symbol map.

\begin{hbdefinition}{def:p02-interpretation-guard}{A finite interpretation guard}
Let \(\alpha:L\to K\) be a \(k\)-term interpretation of a finite
language, with parameter terms \(\bar t\). Put
\[
 D_\alpha(x)=(\mathbf V^\alpha)(x,\bar t).
\]
For a function symbol \(f\in L\) of arity \(n\), including \(n=0\)
for a constant, its closure condition is
\[
 \forall x_1\cdots\forall x_n
 \left(\bigwedge_{j=1}^{n}D_\alpha(x_j)
       \longrightarrow D_\alpha(f^\alpha(\bar x,\bar t))\right).
\]
For \(n=0\) this is the condition that the interpreted constant belongs
to the selected domain. Let \(\operatorname{Inv}_\alpha\) be the finite
conjunction of these conditions. No nonemptiness condition is added:
empty domains without constants are permitted by the source convention.
The guard holds exactly when the indicated pull-back is a structure.
\end{hbdefinition}

\begin{hbtheorem}{thm:p02-guarded-fo-translation}{proved=annex:p2.tr.2}{Guarded first-order translation}
For a first-order sentence \(\varphi\), translate terms by the symbol
interpretation, translate relation atoms likewise, and restrict each
quantifier to \(D_\alpha\). Conjoin the guard for the occurring source
symbols. The resulting sentence holds in a target structure exactly
when its pulled-back source structure is defined and satisfies
\(\varphi\). For universe-preserving zero-term maps, the guard is
vacuous and plain symbol substitution is functorial.
For a sequence of domain-changing operations, retain and apply every
successive domain and guard. This gives the intended sequential truth
condition without asserting equality with the unguarded componentwise
composite of the underlying symbol maps.
\end{hbtheorem}
